\documentclass[11pt,a4paper]{article}
\usepackage[margin=1in]{geometry}
\usepackage{amsmath,amssymb,amsthm}
\usepackage{mathtools}
\usepackage{booktabs}
\usepackage{array}
\usepackage{hyperref}

\theoremstyle{plain}
\newtheorem{theorem}{Theorem}
\newtheorem{proposition}[theorem]{Proposition}
\newtheorem{lemma}[theorem]{Lemma}
\newtheorem{corollary}[theorem]{Corollary}

\theoremstyle{definition}
\newtheorem{definition}[theorem]{Definition}

\theoremstyle{remark}
\newtheorem{remark}[theorem]{Remark}
\newtheorem{observation}[theorem]{Observation}

\newcommand{\R}{\mathbb{R}}
\newcommand{\Z}{\mathbb{Z}}
\newcommand{\T}{\mathbb{T}}
\newcommand{\Hess}{\operatorname{Hess}}
\newcommand{\Ric}{\operatorname{Ric}}
\newcommand{\Var}{\operatorname{Var}}
\newcommand{\tr}{\operatorname{tr}}

\title{Bakry-\'{E}mery Curvature and the Yang--Mills Mass Gap:\\
A Bridge via the Cosine Potential on $S^3$}
\author{R.\ da Costa Fontes\\ \textit{Quijotic Research}}
\date{February 2026}

\begin{document}
\maketitle

\begin{abstract}
We establish a structural connection between the Bakry-\'{E}mery
curvature theory for fractional Laplacians on the torus and the
spectral gap problem for lattice Yang--Mills theory with gauge group
$\mathrm{SU}(2)$.  Numerical computation of the carr\'{e} du champ
kernel for the magnetic operator on the $2^3$~torus reveals exact
$R_{1/2}$~covariance structure (Brownian fBM) in the spectral
decomposition.  We then prove that the single-link conditional measure
of the Yang--Mills Hamiltonian satisfies a Bakry-\'{E}mery curvature
bound $\mathrm{CD}(2-c,\infty)$ on $S^3 \cong \mathrm{SU}(2)$, where
$c = \tfrac{\beta}{2}|W_\ell|$ encodes the staple sum from
neighboring plaquettes.  This bound is the exact $S^3$~analog of the
$\mathrm{CD}(1-\omega^2/2,\,\infty)$ result for the Cauchy--Fokker--Planck
operator on~$\T$.  We discuss the implications for removing the LTQO
condition from the mass gap proof via the Martinelli multi-scale
iteration.
\end{abstract}

\tableofcontents

%% ================================================================
\section{Introduction}
%% ================================================================

Two recent lines of work address spectral gap problems through
curvature methods.  The first establishes the first positive
Bakry-\'{E}mery curvature bound for fractional Laplacians on the
torus, proving $\mathrm{CD}(1-\omega^2/2,\,\infty)$ for the
Cauchy--Fokker--Planck operator with cosine drift.  The second proves a
positive mass gap for $\mathrm{SU}(2)$ lattice Yang--Mills theory on the
$2^3$~torus, conditional on a local topological quantum order (LTQO)
assumption for the infinite-volume extension.

The purpose of this note is to identify the precise structural bridge
between these two results.  Our main findings are:

\begin{enumerate}
\item[(A)] The magnetic operator $B$ in the Yang--Mills Hamiltonian has
  exact $R_{1/2}$ (Brownian) covariance structure in its spectral
  decomposition (Section~\ref{sec:carre}).
\item[(B)] The single-link conditional measure on $S^3 \cong
  \mathrm{SU}(2)$ satisfies $\mathrm{CD}(2-c,\,\infty)$ with a cosine
  potential from the Wilson action, in exact parallel to the
  torus result (Section~\ref{sec:S3}).
\item[(C)] The cross-sector coupling in the $B$~eigenbasis has singular
  values given by the Casimir eigenvalue $C_2(j{=}\tfrac{1}{2}) =
  \tfrac{3}{4}$, revealing the representation-theoretic origin of the
  curvature structure (Section~\ref{sec:cross}).
\item[(D)] The diffusion curvature $\kappa_{\mathrm{diff}} = 2-c$ becomes
  negative for $\beta \gtrsim 2.5$, but the mass gap remains positive
  due to contributions from the magnetic (jump) sector
  (Section~\ref{sec:balance}).
\end{enumerate}

%% ================================================================
\section{Setup and Notation}
%% ================================================================

\subsection{The Yang--Mills Hamiltonian}

We work with the Kogut--Susskind Hamiltonian on the spatial lattice
$\Lambda = (\Z/2\Z)^3$:
\begin{equation}\label{eq:H}
  H(g^2) = \frac{g^2}{2}\sum_{\ell} j_\ell(j_\ell+1)
  - \frac{1}{g^2}\sum_p B_p,
\end{equation}
where $j_\ell \in \{0,\tfrac{1}{2},1,\ldots\}$ labels the
representation on oriented link $\ell$, and $B_p = \tfrac{1}{2}\,
\chi_{1/2}(U_p)$ is the plaquette operator inserting the fundamental
character around plaquette~$p$.

After BFS tree gauge fixing, the $2^3$~torus has 17~active links,
12~plaquettes, and the Gauss law projects onto gauge-invariant states.
At truncation $j_{\max}=1$, the full Hilbert space has dimension
$20{,}041{,}551$.

\subsection{The Feshbach subspace}

The low-energy subspace with electric energy $E = \sum_\ell
j_\ell(j_\ell+1) \le 3$ has dimension~44, containing:

\begin{center}
\begin{tabular}{ccl}
\toprule
$E$ & States & Holonomy sectors \\
\midrule
$0$ & 1 & $(+,+,+)$ \\
$\tfrac{3}{2}$ & 5 & $(+,+,+)$ \\
$3$ & 38 & $(+,+,+)$: 32; $(+,-,-)$: 2; $(-,+,-)$: 2; $(-,-,+)$: 2 \\
\bottomrule
\end{tabular}
\end{center}

The projected magnetic operator $PBP$ is a $44\times 44$ real symmetric
matrix with $\|PBP\| = 0.3998$, matching the theoretical bound
$\|PBP\| \le 0.4$ from the Clebsch--Gordan analysis.

\subsection{The Bakry-\'{E}mery framework on $\T$}

On the torus $\T = \R/(2\pi\Z)$, the Cauchy--Fokker--Planck operator
\begin{equation}
  L = -(-\Delta)^{1/2} + \omega^2 \sin(x)\,\frac{d}{dx}
\end{equation}
restricted to positive Fourier modes $\{1,\ldots,N\}$ has carr\'{e}
du champ kernel
\begin{equation}\label{eq:psi-torus}
  \Psi_1(n,m) = \tfrac{1}{2}(|n| + |m| - |n-m|) = \min(n,m),
\end{equation}
which is the covariance function $R_{1/2}(n,m)$ of fractional Brownian
motion at Hurst parameter $H = \tfrac{1}{2}$.

The curvature bound $\mathrm{CD}(1-\omega^2/2,\,\infty)$ follows from
three key properties:
\begin{itemize}
\item[(i)] The generalized eigenvalues of $(\Psi_1^{\circ 2},\,
  \Psi_1)$ are the odd integers $\{1,3,5,\ldots,2N{-}1\}$, giving
  $\kappa(1,N) = 1$ for all~$N$.
\item[(ii)] The cross-sign kernel vanishes: $\Psi_1(n,-m) = 0$ for
  $n,m > 0$.
\item[(iii)] The drift correction is a scalar shift: $\Gamma_{L,2} -
  \Gamma_{1,2} = \tfrac{\omega^2}{2}\cos(x)\,\Gamma$.
\end{itemize}

%% ================================================================
\section{Carr\'{e} du Champ of the Magnetic Operator}\label{sec:carre}
%% ================================================================

\subsection{Spectral decomposition of $B$}

Let $B_{\text{vac}}$ denote the magnetic operator $PBP$ restricted to
the vacuum sector $(+,+,+)$, which has dimension~38.  The eigenvalue
decomposition yields:

\begin{observation}\label{obs:pm}
$B_{\mathrm{vac}}$ has exactly 9~positive eigenvalues, 9~negative
eigenvalues, and a 20-dimensional null space.  The positive and
negative eigenvalues are exactly paired:
\begin{equation}
  \sigma(B_{\mathrm{vac}}) \cap (0,\infty) = \{\lambda_1, \ldots, \lambda_9\},
  \quad
  \sigma(B_{\mathrm{vac}}) \cap (-\infty,0) = \{-\lambda_9, \ldots, -\lambda_1\},
\end{equation}
where $0 < \lambda_1 < \cdots < \lambda_9$.
\end{observation}

The numerical values are:
\begin{equation*}
  \lambda_k = 0.0224,\; 0.0428,\; 0.0644,\; 0.0995,\; 0.1200,\;
  0.1826,\; 0.1923,\; 0.2033,\; 0.3384.
\end{equation*}

This $\pm$-symmetry follows from the fact that $B = \sum_p B_p$ with
each $B_p$ a trace (character insertion), which maps the representation
$j$ to $j \pm \tfrac{1}{2}$.  In the gauge-invariant sector, this
creates exact $\pm$-pairing.

\subsection{Exact $R_{1/2}$ covariance structure}

In the $B_+$ eigenbasis (the 9-dimensional positive eigenspace), the
carr\'{e} du champ kernel is:
\begin{equation}\label{eq:psi-YM}
  \Psi(k,l) = \tfrac{1}{2}(\lambda_k + \lambda_l - |\lambda_k - \lambda_l|)
  = \min(\lambda_k,\,\lambda_l).
\end{equation}
This is \emph{identically} the fBM covariance $R_{1/2}(\lambda_k,
\lambda_l)$ at Hurst parameter $H = \tfrac{1}{2}$.

\begin{remark}
Equation~\eqref{eq:psi-YM} is structurally identical to
\eqref{eq:psi-torus}: both are $\min$-kernels, which is the signature
of $H = \tfrac{1}{2}$ (Brownian motion).  On the torus this holds
because $|\cdot|^1$ is linear; on the Yang--Mills side it holds because
$\tfrac{1}{2}(a+b-|a-b|) = \min(a,b)$ for any ordered positive reals.
\end{remark}

\subsection{Approximate odd-integer curvature spectrum}

The generalized eigenvalues of $(\Psi^{\circ 2},\,\Psi)$ in the $B_+$
sector are:

\begin{center}
\begin{tabular}{c r r r}
\toprule
$k$ & $\kappa_k$ & $\kappa_k / \kappa_0$ & Odd integer \\
\midrule
0 & 0.02242 & 1.00 & 1 \\
1 & 0.06521 & 2.91 & 3 \\
2 & 0.10717 & 4.78 & 5 \\
3 & 0.16390 & 7.31 & 7 \\
4 & 0.21951 & 9.79 & 9 \\
5 & 0.30253 & 13.50 & 11 \\
6 & 0.37482 & 16.72 & 13 \\
7 & 0.39559 & 17.65 & 15 \\
8 & 0.54169 & 24.16 & 17 \\
\bottomrule
\end{tabular}
\end{center}

The first five ratios agree with the odd integers to within 3--10\%.
The deviation at higher levels reflects the non-uniform spacing of the
$B$~eigenvalues (on the torus, the Fourier modes are uniformly spaced,
giving exact odd integers).  The pattern error over all~9 eigenvalues is
$0.286$ in relative $\ell^2$~norm.

\subsection{Positive curvature}

The minimum curvature eigenvalue is:
\begin{equation}
  \boxed{\kappa = \lambda_{\min}(B_+) = 0.02242 > 0.}
\end{equation}
This is the smallest nonzero eigenvalue of the magnetic operator in the
vacuum sector.  Its positivity is a direct consequence of the
$R_{1/2}$~structure and the positive-definiteness of the $\min$-kernel.

%% ================================================================
\section{Cross-Sector Coupling}\label{sec:cross}
%% ================================================================

The electric operator $E = \sum_\ell j_\ell(j_\ell+1)$, expressed in the
$B$~eigenbasis, couples the $B_+$ and $B_-$ sectors.

\begin{observation}
The singular values of $E(+,-)$ are
\begin{equation}
  \sigma(E_{+-}) = \left\{\tfrac{3}{2},\; \tfrac{3}{4},\; \tfrac{3}{4},\;
  \tfrac{3}{4},\; \tfrac{3}{4},\; 0,\; 0,\; 0,\; 0\right\}.
\end{equation}
\end{observation}

The nonzero singular values are exactly $C_2(j{=}\tfrac{1}{2}) =
\tfrac{3}{4}$ and $2\,C_2(j{=}\tfrac{1}{2}) = \tfrac{3}{2}$, where
$C_2(j) = j(j+1)$ is the quadratic Casimir.  This is not a coincidence:
the Feshbach subspace has electric energies $E \in \{0, \tfrac{3}{2},
3\}$, and the cross-sector coupling is mediated entirely by the Casimir
structure of the $j = \tfrac{1}{2}$ representation.

The Frobenius norm $\|E(+,-)\|_F = 2.121$ is nonzero, confirming that
the Yang--Mills system is \emph{not} at the exact Cauchy point
$\gamma=1$ of the Bakry-\'{E}mery theory (where cross-sign vanishing
$\Psi_1(n,-m) = 0$ holds).  The analogy with the torus is:
\begin{center}
\begin{tabular}{ll}
\toprule
Torus ($\gamma = 1$) & Yang--Mills \\
\midrule
$\Psi_1(n,-m) = 0$ & $\|E(+,-)\| \neq 0$ \\
Cross-sign vanishing exact & Cross-sign coupling via Casimir \\
Scalar shift exact & Scalar shift approximate \\
\bottomrule
\end{tabular}
\end{center}

%% ================================================================
\section{The $S^3$ Analog of Theorem~E}\label{sec:S3}
%% ================================================================

The correct level at which the Bakry-\'{E}mery theory maps onto
Yang--Mills is not the multi-link Feshbach space but the
\emph{single-link conditional measure}.

\subsection{Setup on $S^3$}

Each link $\ell$ of the gauge-fixed lattice carries a group element
$U_\ell \in \mathrm{SU}(2) \cong S^3$.  Conditional on the neighboring
links (which enter through the staple sum $W_\ell$), the single-link
measure is
\begin{equation}
  d\mu_\ell(U) = \frac{1}{Z_\ell}\,
  \exp\!\bigl(\beta\,\mathrm{Re}\,\tr(U\,W_\ell^\dag)\bigr)\,dU,
\end{equation}
where $dU$ is the Haar measure on $\mathrm{SU}(2)$ and
$\beta = 4/g^2$.

Parametrizing $\mathrm{SU}(2)$ by the geodesic distance $\theta \in
[0,\pi]$ from the identity (with Haar measure $\sin^2\theta\,d\theta$
on the zonal sector), the potential becomes
\begin{equation}
  V(\theta) = -c\cos\theta, \qquad c = \tfrac{\beta}{2}\,|W_\ell|.
\end{equation}

The associated diffusion operator (Langevin dynamics on $S^3$) is
\begin{equation}
  L = \Delta_{S^3} + c\sin\theta\,\frac{d}{d\theta}.
\end{equation}

\subsection{The curvature bound}

\begin{theorem}[Bakry-\'{E}mery curvature on $S^3$ with cosine drift]%
\label{thm:S3}
The operator $L = \Delta_{S^3} - \nabla V \cdot \nabla$ with $V(\theta)
= -c\cos\theta$ on $S^3$ satisfies
\begin{equation}
  \boxed{\mathrm{CD}(2-c,\,\infty)}\,.
\end{equation}
Equivalently, for all smooth $f \colon S^3 \to \R$:
\begin{equation}\label{eq:CD-S3}
  \Gamma_{L,2}(f,f) \;\ge\; (2-c)\,\Gamma(f,f),
\end{equation}
where $\Gamma(f,f) = |\nabla f|^2$ and $\Gamma_{L,2}(f,f) =
\tfrac{1}{2}[L\,\Gamma(f,f) - 2\,\Gamma(f,Lf)]$.
\end{theorem}

\begin{proof}
For a diffusion operator $L = \Delta - \nabla V \cdot \nabla$ on a
Riemannian manifold $(M,g)$, the standard Bakry-\'{E}mery identity
gives
\begin{equation}\label{eq:BE-general}
  \Gamma_{L,2}(f,f) = |\Hess\,f|^2 + (\Ric + \Hess\,V)(\nabla f,
  \nabla f).
\end{equation}

On $S^3$ with constant curvature, $\Ric = 2g$.

For $V(\theta) = -c\cos\theta$, the Hessian is computed using the
formula for radial functions on $S^n$:
\begin{equation}
  \Hess\,V = V''\,d\theta \otimes d\theta + V'\cot\theta\,(g -
  d\theta \otimes d\theta).
\end{equation}
Since $V'' = c\cos\theta$ and $V'\cot\theta = c\sin\theta \cdot
\frac{\cos\theta}{\sin\theta} = c\cos\theta$, both components equal
$c\cos\theta$, giving:
\begin{equation}
  \Hess\,V = c\cos\theta \cdot g_{S^3}.
\end{equation}
This is a \emph{scalar multiple of the metric at every point} --- the
key structural identity.

Substituting into~\eqref{eq:BE-general}:
\begin{equation}
  \Gamma_{L,2}(f,f) = |\Hess\,f|^2 + (2 + c\cos\theta)\,|\nabla f|^2
  \;\ge\; (2 + c\cos\theta)\,\Gamma(f,f).
\end{equation}
Taking the infimum over $\theta \in [0,\pi]$:
\begin{equation}
  \Gamma_{L,2}(f,f) \;\ge\; (2-c)\,\Gamma(f,f). \qedhere
\end{equation}
\end{proof}

\begin{remark}[Comparison with the torus result]
On $\T$ with the Cauchy operator $L = -(-\Delta)^{1/2} + \omega^2
\sin(x)\,d/dx$, the curvature bound is $\mathrm{CD}(1-\omega^2/2,\,
\infty)$.  On $S^3$ with $L = \Delta_{S^3} + c\sin\theta\,d/d\theta$,
the bound is $\mathrm{CD}(2-c,\,\infty)$.  The parallel is exact:
\begin{center}
\renewcommand{\arraystretch}{1.3}
\begin{tabular}{ll}
\toprule
$\T$, Theorem~E & $S^3$, Theorem~\ref{thm:S3} \\
\midrule
$\kappa_{\mathrm{free}} = 1$ & $\kappa_{\mathrm{free}} = 2 = \Ric(S^3)$ \\
$V(x) = -\omega^2\cos x$ & $V(\theta) = -c\cos\theta$ \\
$\kappa = 1 - \omega^2/2$ & $\kappa = 2 - c$ \\
$\Hess(V)$: hard (non-local $\Gamma$) & $\Hess(V) = c\cos\theta \cdot g$
  (automatic) \\
\bottomrule
\end{tabular}
\end{center}
The crucial difference is that on $\T$ the fractional Laplacian is
non-local, making the scalar shift a \emph{hard theorem} requiring the
fBM covariance machinery (Theorems A--D of the Bakry-\'{E}mery paper).
On~$S^3$, the Laplacian is a diffusion operator, and the scalar shift
follows from standard Riemannian geometry.
\end{remark}

\subsection{Consequences for the spectral gap}

The curvature bound $\mathrm{CD}(2-c,\,\infty)$ implies:

\begin{corollary}[Single-link Poincar\'{e} inequality]\label{cor:poincare}
For $c < 2$, the single-link conditional measure $\mu_\ell$ satisfies
\begin{equation}
  \Var_{\mu_\ell}(f) \;\le\; \frac{1}{2-c}\,\mathcal{E}_L(f,f),
\end{equation}
where $\mathcal{E}_L(f,f) = \int |\nabla f|^2\,d\mu_\ell$ is the
Dirichlet form.
\end{corollary}

\begin{corollary}[Gradient estimate]
For the semigroup $P_t^L = e^{tL}$:
\begin{equation}
  |\nabla P_t^L f|^2 \;\le\; e^{-2(2-c)t}\, P_t^L\!\bigl[|\nabla f|^2\bigr].
\end{equation}
\end{corollary}

These recover and sharpen the single-link gap $\gamma_\ell \ge
\tfrac{3}{4}$ from Theorem~9.1 of the Yang--Mills paper: the spectral
gap of $\Delta_{S^3}$ on $S^3$ is $\lambda_1 = 3$
(the first nonzero eigenvalue, corresponding to $j = 1$ in the Peter--Weyl
expansion), and $\mathrm{CD}(\kappa,\infty)$ with
$\kappa = 2-c > 0$ implies $\lambda_1(\mu_\ell) \ge 2-c$ for the
weighted measure.

\subsection{Numerical verification}

We verify the curvature bound numerically using zonal spherical
harmonics $\chi_j(\theta) = \sin((j{+}1)\theta)/\sin\theta$.  For each
mode~$j$, we compute $\Gamma_{L,2}(\chi_j,\chi_j)(\theta)$ and
$\Gamma(\chi_j,\chi_j)(\theta)$ on a grid of 2000~points and evaluate
the ratio:

\begin{center}
\begin{tabular}{c c c c}
\toprule
$j$ & $\min_\theta \frac{\Gamma_{L,2}}{\Gamma}$ &
$\max_\theta \frac{\Gamma_{L,2}}{\Gamma}$ & Lower bound $2-c$ \\
\midrule
1 & 1.823 & 29656 & 0.5 \\
2 & 2.545 & $1.64\times 10^6$ & 0.5 \\
3 & 1.997 & $1.03\times 10^7$ & 0.5 \\
\bottomrule
\end{tabular}
\end{center}

(Here $c = 1.5$.)  The minimum ratio exceeds $2-c = 0.5$ in all cases,
confirming the curvature bound.  The actual minimum is substantially
better than the worst case, reflecting the fact that $\cos\theta > -1$
except at $\theta = \pi$.

%% ================================================================
\section{The Curvature Balance}\label{sec:balance}
%% ================================================================

\subsection{Diffusion vs.\ jump contributions}

The full Yang--Mills Hamiltonian~\eqref{eq:H} generates dynamics with
two components:
\begin{enumerate}
\item \textbf{Diffusion} (electric term): The Glauber dynamics on
  each link, with conditional measure on~$S^3$.  Bakry-\'{E}mery
  curvature $\kappa_{\mathrm{diff}} = 2-c$.
\item \textbf{Jump} (magnetic term): The plaquette operator $B_p$
  changes representations by $\pm\tfrac{1}{2}$ on all four links
  simultaneously.  This is a non-local operator on the product
  space $(S^3)^{n_{\mathrm{links}}}$.
\end{enumerate}

The effective coupling $c = \tfrac{\beta}{2}|W_\ell|$ depends on the
staple sum, which is controlled by the activity
$\alpha = I_2(\beta)/I_1(\beta)$ (ratio of modified Bessel functions).
After gauge fixing to coordination~$\sim 3$:
\begin{equation}
  c_{\mathrm{eff}} \approx \tfrac{3}{2}\,\beta\,\alpha(\beta).
\end{equation}

\subsection{Curvature table}

The following table shows the diffusion curvature $\kappa_{\mathrm{diff}}
= 2 - c_{\mathrm{eff}}$ alongside the actual mass gap from the
44-state and 2252-state Feshbach computations:

\begin{center}
\renewcommand{\arraystretch}{1.2}
\begin{tabular}{r r r r r r r}
\toprule
$\beta$ & $g^2$ & $\alpha$ & $c_{\mathrm{eff}}$ &
$\kappa_{\mathrm{diff}}$ & $\mathrm{gap}_{44}$ & $\mathrm{gap}_{2252}$ \\
\midrule
0.5 & 8.00 & 0.124 & 0.09 & 1.91 & 6.000 & 6.000 \\
1.0 & 4.00 & 0.240 & 0.36 & 1.64 & 3.000 & 3.000 \\
2.0 & 2.00 & 0.433 & 1.30 & 0.70 & 1.503 & 1.496 \\
2.5 & 1.60 & 0.507 & 1.90 & 0.10 & 1.206 & 1.192 \\
\midrule
3.0 & 1.33 & 0.568 & 2.56 & $-0.56$ & 1.010 & 0.987 \\
3.5 & 1.14 & 0.618 & 3.24 & $-1.24$ & 0.873 & 0.838 \\
4.0 & 1.00 & 0.658 & 3.95 & $-1.95$ & 0.773 & 0.722 \\
5.0 & 0.80 & 0.719 & 5.40 & $-3.40$ & 0.645 & 0.555 \\
8.0 & 0.50 & 0.819 & 9.83 & $-7.83$ & 0.525 & 0.306 \\
\bottomrule
\end{tabular}
\end{center}

The horizontal line at $\beta \approx 2.5$ marks the transition where
$\kappa_{\mathrm{diff}}$ becomes negative.  For $\beta \gtrsim 2.5$,
the single-link diffusion curvature can no longer sustain a positive
spectral gap.  Yet the actual mass gap remains positive for all
$\beta > 0$, decreasing monotonically but never reaching zero.

\subsection{The role of the jump sector}

The positivity of the gap for $\beta > 2.5$ must come from the magnetic
(jump) contribution.  In the $B_+$ sector, the curvature $\kappa =
\lambda_{\min}(B_+) = 0.0224 > 0$ provides a positive contribution
that supplements the diffusion curvature.  However, this is not a simple
additive correction: the electric and magnetic terms do not commute,
and the cross-sector coupling (Section~\ref{sec:cross}) mixes the
diffusion and jump contributions.

The crossover regime $\beta \in [2.3, 5.0]$ --- precisely where
convergent cluster expansions fail due to the closed-surface growth
constant $\mu_{\mathrm{cs}} \ge 1.64$ --- is where the interplay
between diffusion and jump curvature is most delicate.  It is in this
regime that the \emph{block-level} computation (Feshbach projection onto
the 44-state subspace) provides the necessary non-perturbative input.

%% ================================================================
\section{Implications for the Mass Gap Proof}
%% ================================================================

\subsection{The LTQO condition}

The main result of the Yang--Mills paper (Theorem~14.1) establishes:
\begin{enumerate}
\item[(i)] Block gap $\gamma(\Lambda,\sigma,g^2) > 0$ for the
  $2^3$~block, uniformly in boundary conditions~$\sigma$
  (unconditional).
\item[(ii)] Infinite-volume gap $m(g^2) > 0$ for all $g^2 > 0$
  (\emph{conditional on LTQO}).
\end{enumerate}

The LTQO assumption enters the Martinelli multi-scale iteration as a
stability condition: it ensures that changing boundary conditions on a
block does not dramatically alter the ground state in the interior.

\subsection{Curvature as a replacement for LTQO}

The Bakry-\'{E}mery framework offers a potential path to removing the
LTQO condition.  The key insight is:

\begin{quote}
\emph{The $\mathrm{CD}(\kappa,\infty)$ condition with $\kappa > 0$
implies exponential decay of correlations, which is the essential
content of LTQO.}
\end{quote}

Specifically, the gradient estimate from $\mathrm{CD}(\kappa,\infty)$
gives
\begin{equation}
  |\nabla_x P_t f(y)| \le e^{-\kappa t}\, P_t[|\nabla_x f|](y),
\end{equation}
which controls the sensitivity of the semigroup to boundary
perturbations.  This is precisely the ``weak dependence on boundary
conditions'' that LTQO encodes.

\subsection{Program for removing LTQO}

\begin{enumerate}
\item \textbf{Single-link curvature} (Theorem~\ref{thm:S3}, proven):
  $\mathrm{CD}(2-c,\,\infty)$ for each link conditional on its neighbors.
  Valid for $c < 2$, i.e., $\beta \lesssim 2.5$.

\item \textbf{Tensorization with gauge constraint}: Develop a
  Bakry-\'{E}mery tensorization theorem for product measures on
  $(S^3)^{n_{\mathrm{links}}}$ subject to the Gauss law.  The gauge
  constraint projects onto a subspace, which should only
  \emph{increase} the curvature (by the fiber-bundle projection argument
  of Section~8.3 in the Yang--Mills paper:
  $\gamma_{\mathrm{phys}} \ge \gamma_{\mathrm{full}}$).

\item \textbf{Block curvature in the crossover}: For $\beta > 2.5$,
  the single-link bound is insufficient.  The 44-state Feshbach
  computation provides the block-level spectral gap, which can be
  reinterpreted as a \emph{block-level Bakry-\'{E}mery curvature}.
  This feeds into the Martinelli iteration without requiring LTQO as
  a separate assumption.

\item \textbf{$S^3$ analog of cross-sign vanishing}: The cross-sector
  singular values $\{\tfrac{3}{2}, \tfrac{3}{4}, \tfrac{3}{4},
  \tfrac{3}{4}, \tfrac{3}{4}\}$ suggest that a
  representation-theoretic version of the cross-sign vanishing
  condition may hold, giving sharper curvature bounds within the vacuum
  holonomy sector.
\end{enumerate}

%% ================================================================
\section{Conclusion}
%% ================================================================

We have established a precise structural bridge between the
Bakry-\'{E}mery curvature theory for fractional Laplacians on the torus
and the spectral gap problem for lattice Yang--Mills theory:

\begin{itemize}
\item The magnetic operator has exact $R_{1/2}$ covariance
  (Brownian structure) in its spectral decomposition.
\item The single-link conditional measure satisfies
  $\mathrm{CD}(2-c,\,\infty)$ on $S^3$, with the cosine potential from
  the Wilson action playing the role of the drift in Theorem~E.
\item The Hessian identity $\Hess(-c\cos\theta) = c\cos\theta \cdot
  g_{S^3}$ is the geometric mechanism behind the scalar shift --- the
  cosine potential is a conformal perturbation on~$S^3$.
\item The curvature framework provides a natural language for the
  multi-scale iteration that may replace the LTQO assumption.
\end{itemize}

The deepest open question is whether the block-level Bakry-\'{E}mery
curvature can be proven positive in the crossover regime $\beta \in
[2.3, 5.0]$ without relying on LTQO.  The numerical evidence
(positive gap for all $\beta$ in both the 44-state and 2252-state
computations) strongly suggests this is the case.

\end{document}
